% =====================================================================
%  MODELLO: VERBALE INTERNO  (riunione del solo gruppo)
%
%  1. Copiate questo file nella cartella della milestone in corso:
%       Candidatura/verbali/interni/   RTB/verbali/interni/   PB/verbali/interni/
%     e chiamatelo VI-AAAA-MM-GG.tex (data della riunione), es.
%       RTB/verbali/interni/VI-2026-11-03.tex
%     Dopo il merge la CI mette accanto il PDF: VI-2026-11-03.pdf
%  2. Cambiate i campi qui sotto. Tutto quello che segue il simbolo %
%     è un commento: non viene stampato.
%  3. Scrivete il contenuto dopo \begin{document}.
%  Copertina, registro, presenze, intestazioni e numeri di pagina si
%  fanno da soli. Guida completa: GUIDA.md (ricetta 1).
% =====================================================================
\documentclass[verbale=interno]{gpt}

\titolo{Titolo della riunione}     % di cosa si è parlato, in poche parole
\data{2026-10-26}                  % data della riunione: AAAA-MM-GG
\orario{21:00}{22:00}              % ora di inizio, ora di fine
\luogo{Discord}                    % oppure: Aula 1C150, Torre Archimede
\redattore{guglielmo}              % chi scrive il verbale (la sua chiave)
\verificatore{corro}               % chi lo controlla (un Verificatore,
                                   % diverso dal redattore)

% Chi faceva parte del gruppo il giorno della riunione.
% Lasciate così, a meno che il gruppo sia cambiato.
\componenti{biasini,cescobolla,corro,guglielmo,mazzurana,spagnoli}
\assenti{}                         % chi mancava, es. \assenti{corro,spagnoli}

% Solo DOPO l'assegnazione del capitolato (togliete il % davanti):
% \capitolato{Cn -- Nome del capitolato (Azienda)}

% ---------------------------------------------------------------------
%  REGISTRO DELLE MODIFICHE  --  la riga più recente va IN CIMA.
%  La versione del documento è quella della PRIMA riga.
%
%  \modifica{versione}{data}{autore}{verificatore}{cosa è cambiato}
%
%  Quando il Verificatore approva, aggiunge lui IN CIMA questa riga
%  (con la data del giorno e le chiavi giuste):
%  \modifica{1.0.0}{2026-10-27}{guglielmo}{corro}{Verifica e approvazione}
% ---------------------------------------------------------------------
\modifica{0.1.0}{2026-10-26}{guglielmo}{--}{Prima stesura}

\begin{document}

\section{Ordine del giorno}
\begin{enumerate}
  \item Primo punto.
  \item Secondo punto.
\end{enumerate}

\section{Svolgimento}
\subsection{Primo punto}
Riassunto della discussione.

\subsection{Secondo punto}
Riassunto della discussione.

\section{Decisioni}
\begin{decisioni}
  \decisione{Testo della decisione presa.}
\end{decisioni}

\section{Azioni da intraprendere}
\begin{azioni}
  % \azione[numero issue]{cosa fare}{chi (chiavi, separate da virgole)}{entro quando}
  \azione[\#12]{Descrizione dell'attività.}{guglielmo}{30/10/2026}
\end{azioni}

\section{Prossima riunione}
Data e argomenti previsti.

\end{document}
